\documentclass[11pt,a4paper]{article}
\usepackage[T1]{fontenc}
\usepackage[utf8]{inputenc}
\usepackage{lmodern}
\usepackage[a4paper,top=1.5cm,bottom=1.5cm,left=2.3cm,right=2.3cm]{geometry}
\usepackage{amsmath,amssymb}
\usepackage{xcolor}
\usepackage{longtable,array,booktabs}
\usepackage{fancyhdr}
\newcommand{\barre}[1]{\leavevmode\setbox0\hbox{#1}\rlap{\rule[0.55ex]{\wd0}{0.5pt}}\box0}
\definecolor{bleu}{HTML}{1F4E79}
\pagestyle{fancy}\fancyhf{}
\lhead{\small 5\textsuperscript{e} Mathématiques --- Suites}
\rhead{\small Numérotation à partir de $u_0$}
\cfoot{\small\thepage}
\rfoot{\footnotesize Aide de l'IA (AIA)}
\renewcommand{\headrulewidth}{0.4pt}
\setlength{\parindent}{0pt}
% ---- lignes du tableau ----
\newcommand{\case}{$\square$}
\newcommand{\sect}[1]{\midrule\multicolumn{2}{@{}l}{\rule{0pt}{4ex}\textbf{\large\color{bleu}#1}}\\[0.6ex]\midrule}
\newcommand{\ligne}[1]{#1 & \case\\}                         % savoir-faire simple
\newcommand{\groupe}[1]{\rule{0pt}{3.6ex}\textbf{#1} & \\} % titre d'un groupe
\newcommand{\sous}[1]{\hangindent=2.1em\hangafter=1\hspace{1.2em}$\circ$ #1 & \case\\}  % ligne d'un groupe
\begin{document}
\begin{center}
{\Large\bfseries\color{bleu} Tableau des savoir-faire}\\[1mm]
Les suites numériques
\end{center}
\vspace{-2mm}
\renewcommand{\arraystretch}{1.0}
\begin{longtable}{@{}>{\raggedright\arraybackslash}p{12.6cm}>{\centering\arraybackslash}p{2.9cm}@{}}
\toprule
 & \textbf{Éval. (x ou v)}\\
\endhead
\bottomrule
\endlastfoot

\sect{À revoir avant de commencer}
\ligne{Calculer avec des puissances}
\ligne{Augmenter ou diminuer de $t\,\%$ : multiplier par $1+\frac{t}{100}$ ou par $1-\frac{t}{100}$}
\ligne{Résoudre une équation du 1\textsuperscript{er} degré et un système de deux équations}

\sect{1. Généralités sur les suites}
\groupe{Connaître les définitions}
\sous{l'ensemble des naturels $\mathbb{N}=\{0\,;1\,;2\,;3\,;\ldots\}$, qui numérote les termes à partir de $u_0$}
\sous{une suite : à chaque naturel $n$, on associe un réel $u_n$, le terme de rang $n$}
\sous{une suite définie en extension (liste des termes : $3\,;5\,;7\,;9\,;\ldots$)}
\sous{une suite définie en compréhension (par son terme général $u_n=f(n)$)}
\sous{une suite définie par récurrence ($u_0$ donné et $u_{n+1}=f(u_n)$)}
\sous{le terme général d'une suite}
\sous{une suite croissante, strictement croissante, décroissante, strictement décroissante}
\sous{une suite monotone}
\groupe{Calculer des termes}
\sous{d'une suite définie en fonction du rang : $u_n=f(n)$}
\sous{d'une suite définie par récurrence : $u_0$ donné et $u_{n+1}=f(u_n)$}
\ligne{Passer d'une écriture à l'autre (trouver une formule à partir des premiers termes)}
\ligne{Représenter graphiquement une suite, sur un axe et dans le plan}
\groupe{Déterminer le sens de variation}
\sous{avec le signe de $u_{n+1}-u_n$}
\sous{en comparant $\frac{u_{n+1}}{u_n}$ à $1$ (termes strictement positifs)}
\ligne{Donner un exemple de suite qui vérifie des conditions données}

\sect{2. Suites arithmétiques}
\ligne{Reconnaître une suite arithmétique ($u_{n+1}-u_n$ constant) et donner sa raison $r$}
\ligne{Connaître le terme général : $u_n=u_0+n\,r$, et aussi $u_n=u_p+(n-p)\,r$}
\groupe{Rechercher un terme, la raison ou le rang}
\sous{à partir du premier terme et de la raison}
\sous{à partir de deux termes quelconques}
\groupe{Somme de termes consécutifs}
\sous{calculer $S=\text{nombre de termes}\cdot\frac{\text{premier}+\text{dernier}}{2}$, par exemple $u_0+\ldots+u_{n-1}=n\,\frac{u_0+u_{n-1}}{2}$}
\sous{démontrer cette formule}
\ligne{Décrire la croissance (linéaire) et le sens de variation selon le signe de $r$}

\pagebreak
\sect{3. Suites géométriques}
\ligne{Reconnaître une suite géométrique ($\frac{u_{n+1}}{u_n}$ constant) et donner sa raison $q$}
\ligne{Connaître le terme général : $u_n=u_0\,q^{\,n}$, et aussi $u_n=u_p\,q^{\,n-p}$}
\groupe{Rechercher un terme, la raison ou le rang}
\sous{à partir du premier terme et de la raison}
\sous{à partir de deux termes quelconques}
\sous{avec la moyenne géométrique : $u_n^2=u_{n-1}\,u_{n+1}$}
\groupe{Somme de termes consécutifs}
\sous{calculer $u_0+u_1+\ldots+u_{n-1}=u_0\,\frac{1-q^{\,n}}{1-q}$ \;($q\neq1$)}
\sous{démontrer cette formule}
\ligne{Décrire la croissance (exponentielle) et le sens de variation selon $u_0$ et $q$}
\ligne{Distinguer un accroissement constant ($+\,r$, arithmétique) d'un taux de croissance constant ($\times\, q$, géométrique)}

\sect{4. Limite d'une suite}
\ligne{Définir la limite d'une suite et expliquer cette définition à l'aide d'un schéma}
\ligne{Donner des exemples de suites convergentes et de suites non convergentes}
\ligne{Conjecturer la limite d'une suite à l'aide de la calculatrice ou d'un tableur}
\ligne{Vérifier la valeur d'une limite à l'aide de la définition}
\groupe{Déterminer la limite}
\sous{d'une suite arithmétique (selon le signe de $r$)}
\sous{d'une suite géométrique (selon la valeur de $q$)}
\ligne{Calculer une somme infinie de termes d'une suite géométrique : $u_0+u_1+u_2+\ldots=\frac{u_0}{1-q}$ si $|q|<1$}

\sect{5. Applications financières}
\groupe{Placement à intérêts simples (suite arithmétique)}
\sous{calculer le capital final : $C_n=C_0\,\bigl(1+n\,\frac{t}{100}\bigr)$, avec le taux $t$ en \%}
\sous{trouver le taux, l'intérêt ou la durée}
\groupe{Placement à intérêts composés (suite géométrique)}
\sous{calculer le capital final : $C_n=C_0\,\bigl(1+\frac{t}{100}\bigr)^n$}
\sous{trouver le taux, l'intérêt ou la durée}
\ligne{Comparer intérêts simples et intérêts composés}
\ligne{Calculer la valeur de versements réguliers (annuités) avec une somme géométrique}
\ligne{Réaliser un tableau d'amortissement d'un prêt avec un tableur}

\sect{6. Résoudre un problème}
\ligne{Traduire un énoncé par une suite (récurrence ou terme général) ou par une somme de termes}
\ligne{Résoudre un problème faisant intervenir des suites dans un contexte nouveau (géométrie, population, médicament, \dots)}
\end{longtable}
\end{document}
